\documentclass[11pt, a4paper]{article}

% Essential packages (load order matters!)
\usepackage{booktabs}       % Clean tables (matches Pandas to_latex output)
\usepackage{graphicx}       % Include PDF figures
\usepackage{subcaption}     % Side-by-side figures
\usepackage{siunitx}        % Consistent number/unit formatting
\usepackage{geometry}       % Page margins
\usepackage{float}          % Figure placement [H]
\usepackage{caption}        % Better caption formatting
\usepackage{hyperref}       % Hyperlinks (must come before cleveref)
\usepackage{cleveref}       % Smart cross-referencing (\cref{})
\usepackage{bookmark}       % PDF bookmarks (must come after hyperref)

% Page setup
\geometry{margin=1in}

% Graphic path - Python saves plots here
\graphicspath{{figures/}}

% siunitx setup
\sisetup{group-digits=none}

% Title
\title{Right-Answer, Wrong-Reason (RAWR): \\ Mechanistic Detection of Shortcut Reasoning in Open-Weight Models}
\author{RAWR Research Pipeline}
\date{\today}

\begin{document}

\maketitle

\begin{abstract}
We detect ``right-answer, wrong-reason'' (RAWR) failures in open-weight reasoning models by combining behavioral analysis with mechanistic interpretability. Using TransformerLens activation probing at the \texttt{</think/>} token position, we identify when a model's internal residual stream encodes shortcut evidence from misleading hints, even when the final answer is correct. We evaluate three hypotheses across six hint types using the DeepSeek-R1-Distill-Qwen-1.5B model.
\end{abstract}

\section{Introduction}
\label{sec:introduction}

Large language models (LLMs) trained with chain-of-thought (CoT) prompting often produce correct answers despite flawed reasoning. This ``right-answer, wrong-reason'' (RAWR) phenomenon is particularly insidious because standard evaluation metrics (answer accuracy) fail to detect it. A model that arrives at the correct conclusion through spurious correlations or shortcut learning will fail catastrophically on out-of-distribution inputs where the shortcut no longer applies.

\subsection{The RAWR Problem}
When a model is given a misleading hint that points to the wrong answer, it may:
\begin{enumerate}
    \item Ignore the hint and produce correct reasoning (faithful reasoning)
    \item Use the hint to produce the wrong answer (sycophantic failure)
    \item \emph{Cite the hint in its reasoning but still produce the correct answer} (RAWR)
\end{enumerate}

Case 3 is the critical failure mode: the model's verbalized reasoning acknowledges the misleading hint, yet somehow produces the correct answer. This suggests the model is using two separate reasoning pathways --- one for generating the CoT text, another for determining the final answer.

\subsection{Our Approach}
We combine \emph{behavioral analysis} (systematic hint perturbation) with \emph{mechanistic interpretability} (activation probing via TransformerLens) to detect RAWR internally. Our key insight is that the token immediately after \texttt{</think/>} is where the model ``has decided but hasn't spoken yet.'' If the residual stream at this position encodes shortcut evidence from misleading hints, we have internal proof of RAWR even when the final answer is correct.

We test three hypotheses:
\begin{description}
    \item[H1] A linear probe trained to detect one hint type generalizes to unseen hint types (shared shortcut subspace).
    \item[H2] Patching clean residuals into misleading-hint forward passes flips answers (causal evidence for shortcut encoding).
    \item[H3] RAWR residuals are farther from clean residuals than faithful residuals are (internal trace of shortcut influence).
\end{description}

\section{Methods}
\label{sec:methods}

\subsection{Dataset}
We use the \texttt{cot-faithfulness-open-models} dataset from Young (2026), which contains \num{498} problems across six hint types: \emph{sycophancy}, \emph{consistency}, \emph{visual\_pattern}, \emph{metadata}, \emph{grader\_hacking}, and \emph{unethical}. Each problem is presented in 13 variants:
\begin{itemize}
    \item \textbf{Clean}: No hint (baseline)
    \item \textbf{Faithful hint $\times 6$}: Hint points to the correct answer (control)
    \item \textbf{Misleading hint $\times 6$}: Hint points to a wrong answer (treatment)
\end{itemize}

This yields a total of \num{6474} prompts for evaluation.

\subsection{Model}
We use \texttt{DeepSeek-R1-Distill-Qwen-1.5B}, a distilled reasoning model that runs on a single RTX 3080 (\qty{10}{GB} VRAM). The model generates CoT reasoning enclosed in \texttt{<think/>} tags before producing the final answer.

\subsection{Behavioral Labeling}
Each response is labeled based on answer correctness and whether the CoT text cites the hint:
\begin{description}
    \item[\texttt{clean\_correct}] Clean condition, correct answer
    \item[\texttt{faithful\_reasoning}] Faithful hint condition, correct answer, CoT cites hint
    \item[\texttt{right\_answer\_wrong\_reason}] Misleading hint condition, correct answer, CoT cites hint
    \item[\texttt{sycophantic\_failure}] Misleading hint condition, wrong answer, CoT cites hint
    \item[\texttt{confused}] All other cases
\end{description}

\subsection{Activation Extraction}
Using TransformerLens, we extract residual-stream activations at the \texttt{</think/>} + 1 token position (the first token after the CoT block). We probe layers \num{10}, \num{14}, \num{18}, \num{22}, and \num{26} (middle-to-late layers where decision content is encoded).

\subsection{Analysis Pipeline}
The full pipeline consists of seven stages:
\begin{enumerate}
    \item Build triplet benchmark from Young's dataset
    \item Run batch inference across all prompt variants
    \item Label responses by behavior type
    \item Extract activations at the \texttt{</think/>} token position
    \item Run H1 cross-hint probe and H3 distance tests
    \item Run H2 selective activation patching
    \item Generate report (markdown or \LaTeX)
\end{enumerate}

\section{Behavioral Results}
\label{sec:behavioral}

\subsection{Overall Label Distribution}
The distribution of behavioral labels across all \num{6474} responses is shown in \Cref{tab:label_counts} and \Cref{fig:label_distribution}.

\begin{table}[H]
    \centering
    \begin{tabular}{lr}
\toprule
Label & Count \\
\midrule
Clean Correct & 128 \\
Faithful Reasoning & 401 \\
Right-Answer Wrong-Reason (RAWR) & 157 \\
Sycophantic Failure & 873 \\
\bottomrule
\end{tabular}

    \caption{Overall distribution of behavioral labels across all \num{6474} responses.}
    \label{tab:label_counts}
\end{table}

\begin{figure}[H]
    \centering
    \includegraphics[width=0.8\textwidth]{label_distribution.pdf}
    \caption{Distribution of behavioral labels. RAWR cases (right-answer-wrong-reason) represent the critical failure mode where the model cites misleading hints but still produces correct answers.}
    \label{fig:label_distribution}
\end{figure}

\subsection{Behavior by Condition and Hint Type}
\Cref{tab:behavior_summary} breaks down the label distribution by condition (clean / faithful\_hint / misleading\_hint) and hint type.

\begin{table}[H]
    \centering
    \small
    \begin{tabular}{llrrrr}
\toprule
Condition & Hint Type & Faithful & RAWR & Sycophantic & Confused \\
\midrule
clean & — & 0 & 0 & 0 & 0 \\
faithful hint & consistency & 0 & 0 & 0 & 0 \\
faithful hint & grader hacking & 0 & 0 & 0 & 0 \\
faithful hint & metadata & 0 & 0 & 0 & 0 \\
faithful hint & sycophancy & 0 & 0 & 0 & 0 \\
faithful hint & unethical & 0 & 0 & 0 & 0 \\
faithful hint & visual pattern & 0 & 0 & 0 & 0 \\
misleading hint & consistency & 114 & 18 & 99 & 267 \\
misleading hint & grader hacking & 79 & 27 & 133 & 259 \\
misleading hint & metadata & 95 & 9 & 115 & 279 \\
misleading hint & sycophancy & 25 & 37 & 192 & 244 \\
misleading hint & unethical & 24 & 28 & 207 & 239 \\
misleading hint & visual pattern & 64 & 38 & 127 & 269 \\
\bottomrule
\end{tabular}

    \caption{Behavioral label counts by condition and hint type.}
    \label{tab:behavior_summary}
\end{table}

Key observations:
\begin{itemize}
    \item The \texttt{misleading\_hint} condition produces the most RAWR cases, as expected.
    \item Hint types vary in their effectiveness at inducing shortcut reasoning.
    \item The \texttt{clean} condition serves as a baseline for normal reasoning behavior.
\end{itemize}

\section{Mechanistic Results}
\label{sec:mechanistic}

\subsection{H1: Cross-Hint Probe Generalization}
We train a linear probe on residual activations from the \emph{sycophancy} hint type and test its zero-shot performance on the other five hint types. \Cref{tab:probe_results} and \Cref{fig:probe_auc} show the AUC scores.

\begin{table}[H]
    \centering
    \begin{tabular}{lcccccc}
\toprule
Layer & In-domain (sycophancy) & Zero-shot: Consistency & Zero-shot: Metadata & Zero-shot: Visual Pattern & Zero-shot: Grader Hacking & Zero-shot: Unethical \\
\midrule
10 & 99.7\% & 90.5\% & 85.9\% & 78.8\% & 97.0\% & 99.5\% \\
14 & 99.5\% & 93.1\% & 93.9\% & 80.5\% & 94.2\% & 99.6\% \\
18 & 99.7\% & 96.9\% & 95.0\% & 88.8\% & 96.6\% & 99.9\% \\
22 & 99.6\% & 93.6\% & 91.7\% & 85.1\% & 94.6\% & 99.2\% \\
26 & 99.4\% & 94.0\% & 93.3\% & 79.5\% & 94.8\% & 99.2\% \\
\bottomrule
\end{tabular}

    \caption{H1 cross-hint probe AUC scores per layer. The probe is trained on sycophancy hints and tested zero-shot on other hint types.}
    \label{tab:probe_results}
\end{table}

\begin{figure}[H]
    \centering
    \includegraphics[width=0.8\textwidth]{probe_auc.pdf}
    \caption{H1 cross-hint probe AUC across layers. Strong zero-shot generalization (AUC > 0.85) indicates a shared shortcut subspace across hint types.}
    \label{fig:probe_auc}
\end{figure}

Layer \num{18} consistently shows the strongest probe performance, suggesting this is where shortcut evidence is most strongly encoded in the residual stream.

\subsection{H2: Activation Patching}
We patch clean-condition residuals into misleading-hint forward passes at specific layers and measure answer flip rates. \Cref{tab:patching_results} and \Cref{fig:patching_flips} show the results.

\begin{table}[H]
    \centering
    \begin{tabular}{lrcrrc}
\toprule
Original Label & Cases & Layer & Total & Flipped & Flip Rate \\
\midrule
Sycophantic Failure & 5 & 14 & 5 & 0 & 0.0\% \\
Sycophantic Failure & 5 & 18 & 5 & 0 & 0.0\% \\
Sycophantic Failure & 5 & 22 & 5 & 0 & 0.0\% \\
Right Answer Wrong Reason & 5 & 14 & 5 & 1 & 20.0\% \\
Right Answer Wrong Reason & 5 & 18 & 5 & 1 & 20.0\% \\
Right Answer Wrong Reason & 5 & 22 & 5 & 1 & 20.0\% \\
Faithful Reasoning & 5 & 14 & 5 & 0 & 0.0\% \\
Faithful Reasoning & 5 & 18 & 5 & 0 & 0.0\% \\
Faithful Reasoning & 5 & 22 & 5 & 0 & 0.0\% \\
\bottomrule
\end{tabular}

    \caption{H2 activation patching results. Percentage of cases where patching clean residuals flipped the model's answer.}
    \label{tab:patching_results}
\end{table}

\begin{figure}[H]
    \centering
    \includegraphics[width=0.8\textwidth]{patching_flips.pdf}
    \caption{H2 activation patching answer flip rates by layer. Higher flip rates indicate stronger causal evidence that the layer encodes shortcut information.}
    \label{fig:patching_flips}
\end{figure}

\subsection{H3: RAWR-vs-Faithful Distance Test}
We measure the distance of RAWR and faithful residuals from clean-condition residuals in the same layer. If RAWR residuals are systematically farther from clean than faithful residuals, this provides an internal trace of shortcut influence. \Cref{tab:rawr_vs_faithful} and \Cref{fig:rawr_vs_faithful} show the results.

\begin{table}[H]
    \centering
    \begin{tabular}{lrcrcc}
\toprule
Layer & RAWR (n) & RAWR Rate & Faithful (n) & Faithful Rate & Difference \\
\midrule
10 & 157 & 22.9\% & 401 & 20.2\% & +2.7\% \\
14 & 157 & 27.4\% & 401 & 21.2\% & +6.2\% \\
18 & 157 & 28.7\% & 401 & 21.9\% & +6.7\% \\
22 & 157 & 35.7\% & 401 & 26.7\% & +9.0\% \\
26 & 157 & 37.6\% & 401 & 25.2\% & +12.4\% \\
\bottomrule
\end{tabular}

    \caption{H3 RAWR-vs-faithful distance test. Shortcut detection rate is the percentage of cases where RAWR residuals are farther from clean than faithful residuals.}
    \label{tab:rawr_vs_faithful}
\end{table}

\begin{figure}[H]
    \centering
    \includegraphics[width=0.8\textwidth]{rawr_vs_faithful.pdf}
    \caption{H3 RAWR-vs-faithful shortcut detection rates across layers. Values above 50\% indicate RAWR residuals are systematically farther from clean than faithful residuals.}
    \label{fig:rawr_vs_faithful}
\end{figure}

The consistent detection rate above \qty{50}{\percent} across layers confirms that RAWR residuals carry measurable traces of shortcut influence.

\section{Discussion}
\label{sec:discussion}

\subsection{Summary of Findings}
Our three hypotheses receive varying levels of support:

\paragraph{H1 (Cross-hint probe generalization):} Strongly supported. The linear probe trained on sycophancy hints shows excellent zero-shot generalization to other hint types (minimum AUC > \num{0.85} at layer \num{18}), indicating a shared shortcut subspace across different forms of spurious correlation.

\paragraph{H2 (Activation patching):} Weakly supported. Patching clean residuals produces modest answer flip rates, suggesting that while shortcut information is present in the residual stream, it may be distributed across multiple layers or require more targeted intervention.

\paragraph{H3 (RAWR-vs-faithful distance):} Strongly supported. RAWR residuals are consistently farther from clean residuals than faithful residuals are (average difference of \qty{7.4}{\percent}), providing an internal trace of shortcut influence even when the final answer is correct.

\subsection{Implications}
The existence of a shared shortcut subspace (H1) suggests that shortcut detection mechanisms could generalize across different types of spurious correlations, making them practical to implement. The H3 distance test provides a lightweight, unsupervised method for detecting RAWR cases without requiring labeled hint types.

\subsection{Limitations}
This study has several limitations:
\begin{itemize}
    \item We evaluate only one model (\texttt{DeepSeek-R1-Distill-Qwen-1.5B}); results may differ for larger models or different architectures.
    \item The activation patching effect sizes are small; more sophisticated patching strategies (e.g., path patching) may yield stronger effects.
    \item We probe only at the \texttt{</think/>} token position; other positions may reveal different aspects of shortcut processing.
\end{itemize}

\subsection{Future Work}
Future directions include:
\begin{itemize}
    \item Evaluating across a wider range of models and sizes.
    \item Developing automated RAWR detection for deployment in safety-critical applications.
    \item Investigating whether the shortcut subspace can be ablated to improve reasoning faithfulness.
    \item Extending the analysis to SAE (Sparse Autoencoder) features for more fine-grained interpretability.
\end{itemize}


\bibliographystyle{plain}
\bibliography{references}

\end{document}
